% Generated draft: requires mathematical and source review.
\documentclass[11pt]{article}
\usepackage[margin=25mm]{geometry}
\usepackage{fontspec}
\setmainfont{Latin Modern Roman}
\usepackage{amsmath,amssymb,mathtools,unicode-math}
\setmathfont{Latin Modern Math}
\usepackage{enumitem,needspace,xurl,hyperref}
\setlength{\parindent}{0pt}
\setlength{\parskip}{5pt}
\setlength{\emergencystretch}{3em}
\newcommand{\N}{\mathbb{N}}
\newcommand{\Z}{\mathbb{Z}}
\newcommand{\Q}{\mathbb{Q}}
\newcommand{\R}{\mathbb{R}}
\newcommand{\C}{\mathbb{C}}
\newcommand{\F}{\mathbb{F}}
\newcommand{\A}{\mathbb{A}}
\newcommand{\PP}{\mathbb{P}}
\newcommand{\E}{\mathbb{E}}
\newcommand{\eps}{\varepsilon}
\newcommand{\dd}{\,\mathrm{d}}
\begin{document}
\section*{UnsolvedMath ID: 701236}
\textbf{Source problem code:} catalogue entry
\subsection{Definitions and mathematical statement}
For a group \(G\) with a topology \(\tau\), being a topological group means that multiplication \(G\times G\to G\) and inversion \(G\to G\) are continuous, where the domain of multiplication has the product topology. The topology is non-discrete if not every singleton is open. A subset \(D\subseteq G\) is discrete if, for every \(d\in D\), there is an open set \(U\subseteq G\) such that \(U\cap D=\{d\}\). Define its set of accumulation points in the ambient group by
\[
D'_G=\{x\in G: \text{every open neighbourhood }U\text{ of }x\text{ satisfies }U\cap(D\setminus\{x\})\ne\varnothing\}.
\]
The existence question attributed to E. G. Zelenyuk is whether ZFC proves that there exists a countable non-discrete topological group \((G,\tau)\) such that
\[
\forall D\subseteq G,\qquad D\text{ discrete}\ \Longrightarrow\ \neg\exists p\in G\ (D'_G=\{p\}).
\]
Here countable means \(|G|\leq\aleph_0\). The supplied record does not specify a separation axiom for topological groups. If the intended convention requires Hausdorffness, that hypothesis must be added; in that case countability and non-discreteness force \(G\) to be countably infinite. The intended separation convention remains unresolved. Accumulation points in the displayed formulation are counted throughout \(G\), including outside \(D\).

The source additionally reports that such a group is known to exist under Martin's Axiom. This is retained as a supplied qualification, without independent verification.
\subsection{Short English statement}
Does ZFC prove the existence of a countable non-discrete topological group having no discrete subset with exactly one accumulation point in the group? The source reports existence under Martin's Axiom. The intended separation axiom for the group topology requires clarification.
\subsection{Sources}
\begin{itemize}
\item Supplied local UnsolvedMath registry record: permanent ID \(701236\), exact source problem code \texttt{catalogue entry}, title \emph{Kourovka Notebook Problem 16.44}. The registry labels the problem open; this is supplied metadata, not an independently verified status assessment.
\item \emph{The Kourovka Notebook (Unsolved Problems in Group Theory)}, original problem 16.44, 16th issue (2006); attributed to E. G. Zelenyuk. The supplied background also identifies the 21st issue edition (2026), arXiv:1401.0300v48. The collection metadata names E. I. Khukhro and V. D. Mazurov as editors.
\item The supplied statement reports existence under Martin's Axiom, but supplies no separate reference for that result. No external source was consulted.
\end{itemize}
\end{document}
